\section{Conclusion}
\sys{} evaluates whether LLM agents can move from a CVE identifier to auditable evidence of IoT firmware vulnerability reproduction. By combining a manually curated 30-CVE public-evidence dataset, 450 reproduction runs, 450 scoring passes, and post-hoc human analysis, it measures progress that binary metrics would obscure. The benchmark is designed to reveal where agents lose grounding: public vulnerability understanding, artifact acquisition, firmware extraction, binary identification, rehosting, or real-target triggering. Its central lesson is that future cyber agents must be evaluated not only on plausible reports or PoC text, but on whether their actions remain grounded in the real artifacts they claim to reproduce.
